% GENERATED FILE. Edit canonical lessons, then run npm run generate:books.
% canonical-source-hash: fnv1a64:1043180c87c6768d
% canonical-lessons: FR-C423-croisiere, FR-C423-accoster, FR-C423-tanguer, FR-C423-chavirer, FR-C423-deriver, FR-R423-revision-416-423

\chapter{En croisière}
\label{ch:fr-en-croisiere}
\hlchaptermodality{423}

\begin{chapteropening}
\textbf{By the end of this chapter:} \emph{I can talk about a cruise: the cruise, coming alongside, pitching, capsizing and drifting.}

Complete the last lesson of chapter 423: I can talk about a cruise: the cruise, coming alongside, pitching, capsizing and drifting.
\end{chapteropening}

\section[la croisière]{la croisière — the cruise}

\label{lesson:FR-C423-croisiere}



\begin{quote}
Before the new one: say the French for the waistcoat, then the French for to get off.
\end{quote}

\subsection*{You'll want to know: la croisière}
\textbf{la croisière} — \textquotedblleft{}the cruise\textquotedblright{}.

\begin{quote}
Mes grands-parents ont fait une croisière sur le fleuve.
\end{quote}

\begin{quote}
\emph{My grandparents went on a river cruise.}
\end{quote}

\emph{Une vitesse de croisière} — cruising speed.

\begin{cousinweb}
\emph{croisière} is built on \emph{croiser} 'to cross, to sail to and fro', from \emph{croix} 'cross', which is inherited from Latin \emph{crux} (Lewis \& Short).
\end{cousinweb}

\begin{grammarlens}[title={What you've built}]
One more word for telling, explaining and describing at B1.
\end{grammarlens}

\subsection*{Your turn}
\begin{itemize}
\raggedright
  \item \emph{Say it:} \emph{la croisière}
  \item \emph{Say it:} \emph{la croisière}, once more
  \item \emph{Say it:} a sentence of your own with \emph{la croisière}
  \item \emph{Recall:} say the French for the waistcoat, then the French for to get off, then say \emph{la croisière} again
\end{itemize}

\begin{tcolorbox}[breakable,colback=teal!4,colframe=teal!35!black,title={Before you move on}]
What does la croisière mean? (\textquotedblleft{}The cruise\textquotedblright{}.) Say it once more.
\end{tcolorbox}
\section[accoster]{accoster — to come alongside, to dock; to approach (someone)}

\label{lesson:FR-C423-accoster}



\begin{quote}
Before the new one: say the French for to get off, then the French for the cruise.
\end{quote}

\subsection*{You'll want to know: accoster}
\textbf{accoster} — \textquotedblleft{}to come alongside, to dock; to approach (someone)\textquotedblright{}.

\begin{quote}
Le bateau accoste au port vers midi.
\end{quote}

\begin{quote}
\emph{The boat docks in the harbour around noon.}
\end{quote}

\emph{Accoster quelqu'un} — to go up to someone in the street.

\begin{cousinweb}
\emph{accoster} is built on \emph{coste}, the old form of \emph{côte} 'side, coast', which is inherited from Latin \emph{costa} 'rib, side' (Lewis \& Short): to bring the ship's side to the shore.
\end{cousinweb}

\begin{grammarlens}[title={What you've built}]
One more word for telling, explaining and describing at B1.
\end{grammarlens}

\subsection*{Your turn}
\begin{itemize}
\raggedright
  \item \emph{Say it:} \emph{accoster}
  \item \emph{Say it:} \emph{accoster}, once more
  \item \emph{Say it:} a sentence of your own with \emph{accoster}
  \item \emph{Recall:} say the French for to get off, then the French for the cruise, then say \emph{accoster} again
\end{itemize}

\begin{tcolorbox}[breakable,colback=teal!4,colframe=teal!35!black,title={Before you move on}]
What does accoster mean? (\textquotedblleft{}To come alongside, to dock; to approach (someone)\textquotedblright{}.) Say it once more.
\end{tcolorbox}
\section[tanguer]{tanguer — to pitch (of a boat); to sway}

\label{lesson:FR-C423-tanguer}



\begin{quote}
Before the new one: say the French for the cruise, then the French for to come alongside.
\end{quote}

\subsection*{You'll want to know: tanguer}
\textbf{tanguer} — \textquotedblleft{}to pitch (of a boat); to sway\textquotedblright{}.

\begin{quote}
Avec le vent, le bateau tangue beaucoup.
\end{quote}

\begin{quote}
\emph{In the wind the boat pitches a lot.}
\end{quote}

\emph{Tanguer} is front to back; \emph{rouler} is side to side.

\begin{cousinweb}
The origin of \emph{tanguer} is not settled. It is a sailors' word, and EtymDB gives no reliable source for it.
\end{cousinweb}

\begin{grammarlens}[title={What you've built}]
One more word for telling, explaining and describing at B1.
\end{grammarlens}

\subsection*{Your turn}
\begin{itemize}
\raggedright
  \item \emph{Say it:} \emph{tanguer}
  \item \emph{Say it:} \emph{tanguer}, once more
  \item \emph{Say it:} a sentence of your own with \emph{tanguer}
  \item \emph{Recall:} say the French for the cruise, then the French for to come alongside, then say \emph{tanguer} again
\end{itemize}

\begin{tcolorbox}[breakable,colback=teal!4,colframe=teal!35!black,title={Before you move on}]
What does tanguer mean? (\textquotedblleft{}To pitch (of a boat); to sway\textquotedblright{}.) Say it once more.
\end{tcolorbox}
\section[chavirer]{chavirer — to capsize; to tip over}

\label{lesson:FR-C423-chavirer}



\begin{quote}
Before the new one: say the French for to come alongside, then the French for to pitch.
\end{quote}

\subsection*{You'll want to know: chavirer}
\textbf{chavirer} — \textquotedblleft{}to capsize; to tip over\textquotedblright{}.

\begin{quote}
La petite barque a chaviré près des rochers.
\end{quote}

\begin{quote}
\emph{The little boat capsized near the rocks.}
\end{quote}

Figurative: \emph{son cœur chavire} — her heart turns over.

\begin{cousinweb}
\emph{chavirer} was borrowed from Provençal \emph{capvirar} 'to turn upside down' (EtymDB), from \emph{cap} 'head', from Latin \emph{caput}, and \emph{virar} 'to turn'.
\end{cousinweb}

\begin{grammarlens}[title={What you've built}]
One more word for telling, explaining and describing at B1.
\end{grammarlens}

\subsection*{Your turn}
\begin{itemize}
\raggedright
  \item \emph{Say it:} \emph{chavirer}
  \item \emph{Say it:} \emph{chavirer}, once more
  \item \emph{Say it:} a sentence of your own with \emph{chavirer}
  \item \emph{Recall:} say the French for to come alongside, then the French for to pitch, then say \emph{chavirer} again
\end{itemize}

\begin{tcolorbox}[breakable,colback=teal!4,colframe=teal!35!black,title={Before you move on}]
What does chavirer mean? (\textquotedblleft{}To capsize; to tip over\textquotedblright{}.) Say it once more.
\end{tcolorbox}
\section[dériver]{dériver — to drift; to derive}

\label{lesson:FR-C423-deriver}



\begin{quote}
Before the new one: say the French for to pitch, then the French for to capsize.
\end{quote}

\subsection*{You'll want to know: dériver}
\textbf{dériver} — \textquotedblleft{}to drift; to derive\textquotedblright{}.

\begin{quote}
Sans moteur, le canot a dérivé toute la nuit.
\end{quote}

\begin{quote}
\emph{Without an engine the dinghy drifted all night.}
\end{quote}

\emph{À la dérive} — adrift.

\begin{cousinweb}
\emph{dériver} 'to drift' is usually linked to English \emph{drive}, reshaped after the older \emph{dériver} 'to divert a stream', from Latin \emph{dērīvāre} (Lewis \& Short). The details are not settled.
\end{cousinweb}

\begin{grammarlens}[title={What you've built}]
One more word for telling, explaining and describing at B1.
\end{grammarlens}

\subsection*{Your turn}
\begin{itemize}
\raggedright
  \item \emph{Say it:} \emph{dériver}
  \item \emph{Say it:} \emph{dériver}, once more
  \item \emph{Say it:} a sentence of your own with \emph{dériver}
  \item \emph{Recall:} say the French for to pitch, then the French for to capsize, then say \emph{dériver} again
\end{itemize}

\begin{tcolorbox}[breakable,colback=teal!4,colframe=teal!35!black,title={Before you move on}]
What does dériver mean? (\textquotedblleft{}To drift; to derive\textquotedblright{}.) Say it once more.
\end{tcolorbox}
\section[(dialogue)]{Retour sur les chapitres 416–423}

\label{lesson:FR-R423-revision-416-423}



\begin{quote}
Say the French for the flight and the airport building.
\end{quote}

\subsection*{Your turn}
Say these aloud:

\begin{itemize}
\raggedright
  \item the flight, the airport building, the check-in, the hold, the boarding
  \item the customs officer, to declare, the scanner gate, to search, to change planes
  \item to board, the flight attendant, the steward, the porthole, the aisle
  \item to take off, the take-off, the landing, to land, to fly over
  \item the turbulence, the stopover, to shift, the gap, to delay
  \item the cancellation, the refund, to confirm, to cancel, the overbooking
  \item the ferry, the crossing, the cabin, the waistcoat, to get off
  \item the cruise, to come alongside, to pitch, to capsize, to drift
\end{itemize}

\begin{tcolorbox}[breakable,colback=teal!4,colframe=teal!35!black,title={Before you move on}]
The flight? (\textbf{le vol}.) The airport building? (\textbf{l'aérogare}.) The check-in? (\textbf{l'enregistrement}.) The hold? (\textbf{la soute}.) The boarding? (\textbf{l'embarquement}.) The customs officer? (\textbf{le douanier}.) To declare? (\textbf{déclarer}.) The scanner gate? (\textbf{le portique}.) To search? (\textbf{fouiller}.) To change planes? (\textbf{transiter}.) To board? (\textbf{embarquer}.) The flight attendant? (\textbf{l'hôtesse}.) The steward? (\textbf{le steward}.) The porthole? (\textbf{le hublot}.) The aisle? (\textbf{l'allée}.) To take off? (\textbf{décoller}.) The take-off? (\textbf{le décollage}.) The landing? (\textbf{l'atterrissage}.) To land? (\textbf{atterrir}.) To fly over? (\textbf{survoler}.) The turbulence? (\textbf{la turbulence}.) The stopover? (\textbf{l'escale}.) To shift? (\textbf{décaler}.) The gap? (\textbf{le décalage}.) To delay? (\textbf{retarder}.) The cancellation? (\textbf{l'annulation}.) The refund? (\textbf{le remboursement}.) To confirm? (\textbf{confirmer}.) To cancel? (\textbf{décommander}.) The overbooking? (\textbf{la surréservation}.) The ferry? (\textbf{le ferry}.) The crossing? (\textbf{la traversée}.) The cabin? (\textbf{la cabine}.) The waistcoat? (\textbf{le gilet}.) To get off? (\textbf{débarquer}.) The cruise? (\textbf{la croisière}.) To come alongside? (\textbf{accoster}.) To pitch? (\textbf{tanguer}.) To capsize? (\textbf{chavirer}.) To drift? (\textbf{dériver}.)
\end{tcolorbox}
